
\vspace{2em}
\begin{center}
\resizebox{0.9\textwidth}{!}{
\begin{tikzpicture}[
    line width=1.2pt,
    line cap=round,
    line join=round
]
% 1. Wheatstone Bridge
\draw (-2,0) to[R] (0,2);
\draw (0,2) to[R] (2,0);
\draw (2,0) to[R] (0,-2);
\draw (0,-2) to[R] (-2,0);
\node at (-1.5, 1.2) {$R$};
\node at (1.8, 1.2) {$R + \Delta R$};
\node at (-2.1, -1.2) {$R + \Delta R$};
\node at (1.5, -1.2) {$R$};
\draw (0,2) -- (0,2.7);
\draw (-0.5, 2.7) -- (0.5, 2.7);
\node at (0, 3.1) {$V_{DD}$};
\fill (0,2) circle (2.5pt); 
\draw (0,-2) -- (0,-2.7);
\draw (-0.4, -2.7) -- (0.4, -2.7);
\draw (-0.25, -2.85) -- (0.25, -2.85);
\draw (-0.1, -3.0) -- (0.1, -3.0);
\fill (0,-2) circle (2.5pt); 

% 2. Instrumentation Amplifier
\draw (4.5, 0.8) rectangle (5.2, -1.8);
\draw (5.2, 0.8) -- (7.8, -0.5) -- (5.2, -1.8);
\node at (4.85, 0) {$\boldsymbol{-}$};
\node at (4.85, -1) {$\boldsymbol{+}$};

% 3. Lowpass filter
\draw (9.5, 1) rectangle (13.5, -2);
\draw (10.3, 0.5) -- (11.8, 0.5) to[out=0, in=100] (12.6, -1.2);

% 4. ADC
\draw (14.8, -0.5) -- (15.5, 0.8) -- (17.3, 0.8) -- (17.3, -1.8) -- (15.5, -1.8) -- cycle;
\node[font=\sffamily] at (16.4, -0.5) {ADC};

% 5. Routing
\draw (2,0) -- (4.5,0);
\draw (-2,0) -- (-3.2,0) -- (-3.2,-2.4) -- (3.2,-2.4) -- (3.2,-1) -- (4.5,-1);
\node at (3.8, -0.5) {$V_{in}$};
\draw (7.8, -0.5) -- (9.5, -0.5);
\draw (13.5, -0.5) -- (14.8, -0.5);
\draw (17.3, -0.5) -- (19, -0.5);
\draw[fill=white] (19.1, -0.5) circle (0.1);
\draw (18.0, -0.75) -- (18.3, -0.25);
\node at (19.8, -0.5) {$D_{out}$};
\end{tikzpicture}
}
\end{center}

\vspace{0.5em}
\begin{center}
\textbf{Figure 1. Strain gage signal condition architecture.}
\end{center}
\vspace{1em}

\textbf{Fig 1.} show the sensor signal conditioning ``front-end'' architecture for a strain gage sensor. The differential output voltage of the strain gage is amplified by the instrumentation amplifier and converted to a single-ended voltage. The output of the amplifier is fed to a lowpass filter which sets the noise 